% =====================================================================================
%  Thai business groups: paper template in arXiv preprint style.
%  Every number comes from generated/macros.tex and generated/tab_*.tex, written by
%  `python -m tbg paper`. Edit THIS file (repo/paper/paper.tex); the workspace copy is
%  overwritten on each build. `python -m tbg paper` compiles with pdfLaTeX (the arXiv
%  default) and switches to XeLaTeX only when Thai script appears in names.
% =====================================================================================
\documentclass[11pt]{article}
\usepackage{iftex}
\newif\ifunicodeengine
\ifXeTeX\unicodeenginetrue\fi
\ifLuaTeX\unicodeenginetrue\fi
\usepackage{amsmath}
\ifunicodeengine
  \usepackage{fontspec}
  \IfFontExistsTF{texgyretermes-regular.otf}{\setmainfont{texgyretermes}[Extension=.otf,
    UprightFont=*-regular, BoldFont=*-bold, ItalicFont=*-italic, BoldItalicFont=*-bolditalic]}{}
  \IfFileExists{unicode-math.sty}{\usepackage{unicode-math}%
    \IfFontExistsTF{texgyretermes-math.otf}{\setmathfont{texgyretermes-math.otf}}{}}{}
\else
  \usepackage[T1]{fontenc}
  \IfFileExists{newtxtext.sty}{\usepackage{newtxtext}\usepackage{newtxmath}}{\usepackage{mathptmx}}
\fi
\usepackage[letterpaper,textwidth=6.5in,textheight=9in,top=1in,headheight=14pt,headsep=25pt,footskip=30pt]{geometry}
\usepackage{graphicx}
\usepackage{booktabs,tabularx}
\usepackage[round,authoryear]{natbib}
\usepackage{xcolor}
\usepackage[font=small,labelfont=bf]{caption}
\usepackage{microtype}
\usepackage{fancyhdr}
\usepackage{hyperref}
\setlength{\parindent}{0pt}
\setlength{\parskip}{5.5pt}

% ---- Thai names: XeLaTeX/LuaLaTeX only (pdfLaTeX builds never contain Thai script) ----
\newif\ifthaifound  % defined outside the conditional so pdfLaTeX can skip the block safely
\ifunicodeengine
  \newcommand{\trythaifont}[1]{\ifthaifound\else\IfFontExistsTF{#1}{\newfontfamily\thaifont[Script=Thai]{#1}\global\thaifoundtrue}{}\fi}
  \trythaifont{Sarabun}\trythaifont{TH Sarabun New}\trythaifont{Noto Serif Thai}\trythaifont{Noto Sans Thai}
  \trythaifont{Leelawadee UI}\trythaifont{Tahoma}\trythaifont{Loma}\trythaifont{Garuda}\trythaifont{Norasi}
  \ifthaifound\else\let\thaifont\rmfamily\fi
  \ifXeTeX\newcommand{\thai}[1]{{\thaifont\XeTeXlinebreaklocale "th"\XeTeXlinebreakskip=0pt plus 1pt #1}}
  \else\newcommand{\thai}[1]{{\thaifont #1}}\fi
\else
  \newcommand{\thai}[1]{#1}
\fi

\newif\ifsynthetic
\InputIfFileExists{generated/macros.tex}{}{\errmessage{generated/macros.tex is missing: run python -m tbg paper}}
\hypersetup{colorlinks=true, linkcolor=blue!50!black, citecolor=blue!50!black, urlcolor=blue!60!black,
  pdftitle={\PaperTitle}, pdfauthor={\PaperAuthors}, pdfkeywords={\PaperKeywordsPlain}}

% ---- arXiv preprint look: running head, ruled title block, compact section styles ----
\pagestyle{fancy}
\fancyhf{}
\fancyhead[L]{\small\scshape \PaperShortTitle}
\fancyhead[R]{\small\scshape A Preprint}
\fancyfoot[C]{\thepage}
\renewcommand{\headrulewidth}{0.4pt}
\makeatletter
\renewcommand\section{\@startsection{section}{1}{\z@}{-2.0ex \@plus -0.5ex \@minus -0.2ex}%
  {1.5ex \@plus 0.3ex \@minus 0.2ex}{\large\bfseries\raggedright}}
\renewcommand\subsection{\@startsection{subsection}{2}{\z@}{-1.8ex \@plus -0.5ex \@minus -0.2ex}%
  {0.8ex \@plus 0.2ex}{\normalsize\bfseries\raggedright}}
\renewcommand\paragraph{\@startsection{paragraph}{4}{\z@}{1.5ex \@plus 0.5ex \@minus 0.2ex}{-1em}%
  {\normalsize\bfseries}}
\makeatother
\renewcommand{\maketitle}{%
  \thispagestyle{plain}%
  \par\vbox{\hrule height 4pt\vskip 0.25in
    {\centering\LARGE\bfseries \PaperTitle\par}%
    \vskip 0.22in\hrule height 1pt}%
  \vskip 0.12in
  {\centering{\scshape A Preprint}\par
   \vskip 0.25in\PaperAuthorBlock\par
   \vskip 0.2in\today\par}%
  \vskip 0.15in}
\renewenvironment{abstract}{\centerline{\large\bfseries Abstract}\vspace{0.3ex}\begin{quote}}{\end{quote}}
\newcommand{\keywords}[1]{\par\noindent\textbf{\textit{Keywords}}\enspace #1\par\medskip}

\begin{document}
\maketitle

\begin{abstract}
\noindent Thailand's corporate sector is organised around business groups: families, the state and a few foreign
parents that control many firms through chains of ownership. No open, firm-level map of these groups exists. We
build one from public sources only: shareholder registers published by the Stock Exchange of Thailand, subsidiary
disclosures in the Form 56-1 One Report, shareholder lists filed with the Department of Business Development, bank
Pillar~3 reports and the Global LEI System. The network links \NEntities{} entities through \NHoldings{} ownership
ties. Tracing every chain to its end with weakest-link control rights, we identify \NGroupsMulti{} groups that control
two or more firms and \NGroupsTenPlus{} that control ten or more. At the \ThresholdPct{} control threshold, families
are the ultimate owners of \ShareFamilyTwenty{} of listed firms and the state of \ShareStateTwenty{};
\ShareWidelyTwenty{} are widely held. Among controlled listed firms, \SharePyramid{} sit inside pyramids, and
control rights exceed cash-flow rights by \MeanWedge{} percentage points on average. Bank groups share directors with
\ShareBigGroupsBankInterlock{} of business groups that have ten or more firms. Code and data pipeline are released so
the map can be rebuilt every year.
\ifsynthetic\par\smallskip\textit{Remark.} This version is a demonstration built from synthetic data: all companies, persons and figures in it are fictional.\fi
\end{abstract}
\PaperKeywordsLine

\section{Introduction}
Business groups dominate private economic activity in many emerging economies \citep{khanna2007,morck2005}. In
Thailand, a small number of family conglomerates, state enterprises and bank-centred groups have shaped capital
accumulation since the mid-twentieth century \citep{suehiro1989}, and the 1997 crisis restructured rather than dissolved
them \citep{polsiri2006,pasuk2008}. Who owns whom therefore matters for questions as different as competition in
retail and telecommunications, the protection of minority shareholders, and the exposure of banks to a handful of
borrowers.

Measuring groups is hard because a group is not a legal entity. It is defined by control relationships that run
through listed and unlisted companies, Thai and foreign registries, and nominee or depositary accounts. Existing
evidence for Thailand relies on samples of listed firms and their largest shareholders
\citep{claessens2000,wiwattanakantang2001}. The unlisted layer, where most subsidiaries and the holding companies at
the top of each group sit, remains largely unmapped.

This paper builds an open and reproducible map of that layer and uses it to answer three questions.
\textbf{RQ1:} Who ultimately controls Thai listed firms and the subsidiaries beneath them, and how many groups are
there? \textbf{RQ2:} How far do control rights diverge from cash-flow rights, and how deep are the pyramids?
\textbf{RQ3:} How are financial institutions connected to business groups through ownership, shared directors and
fund holdings?

We make three contributions. First, coverage: the network includes every subsidiary, associate and joint venture
that listed firms are required to disclose, plus the private holding companies above them. Second, measurement: a
two-stage ultimate-owner algorithm handles cross-holdings and cycles, aggregates family members, and excludes
depositary, custodian and fund accounts that hold shares without controlling them. Third, a description of the
finance--business nexus that combines three kinds of links. Section~\ref{sec:context} describes the disclosure
regime, Section~\ref{sec:data} the data, Section~\ref{sec:method} the measures, Section~\ref{sec:results} the
results and Section~\ref{sec:discussion} the limitations.

\section{Institutional background}\label{sec:context}
\paragraph{Listed companies.} Since fiscal year 2021, listed companies file an annual registration statement and
annual report combined as the Form 56-1 One Report. It contains a chart and a list of the group's shareholding
structure, and the notes to the financial statements report investments in subsidiaries, associates and joint
ventures with direct and effective ownership percentages. The Stock Exchange of Thailand (SET) publishes the ten
largest shareholders of each company at the latest book-closing date, and holders must report to the Securities and
Exchange Commission (SEC) whenever their stake crosses a multiple of 5\% (Form 246-2).

\paragraph{Depositary receipts and nominees.} Thai NVDR Co., Ltd., a subsidiary of the SET, issues non-voting
depositary receipts and therefore appears among the largest shareholders of many companies without exercising
control. Foreign custodians and nominee accounts also appear under their own names. Treating either as an owner would
overstate dispersed and foreign ownership, so we exclude both from the set of possible controllers.

\paragraph{Private companies.} Private limited companies file a shareholder list (form BOJ.5) and audited financial
statements with the Department of Business Development (DBD) every year. Each company has a 13-digit registration
number that also serves as its tax identification number, which we use to link records across sources.

\paragraph{Financial groups.} The Bank of Thailand supervises banks on a consolidated basis as financial business
groups, and banks typically disclose the composition of their group in Pillar~3 reports. Several groups now operate
through a holding company above the bank (for example SCB~X since 2022), so the bank is not always the apex.

\paragraph{Foreign ownership.} Under the Foreign Business Act B.E.~2542 (1999), a company in which foreigners hold half
or more of the shares is treated as foreign for restricted businesses. Some firms therefore combine a Thai majority
of shares with voting arrangements or preference shares that give control elsewhere. Share percentages are thus a
lower bound on the complexity of control, a point we return to in Section~\ref{sec:discussion}.

\section{Data}\label{sec:data}
Table~\ref{tab:sources} lists the sources and Table~\ref{tab:coverage} reports coverage. Shareholdings are as of \AsOfDate{} and subsidiary lists come from fiscal year \DataYear{} reports. Subsidiary tables were
extracted from the One Report PDFs with a large language model instructed to copy, not infer, every row; low-confidence
rows and a random share of documents were checked by hand. Names were normalised (legal forms, titles, Thai PDF glyph
artefacts) and matched with exact keys, registration numbers and a conservative fuzzy rule whose borderline cases were
reviewed manually. When two sources report the same link we keep, in order, DBD shareholder lists, SET registers and
direct holdings in the One Report; effective holdings are used only when no direct figure is available.

\begin{table}[t]
\centering\small
\caption{Public data sources}\label{tab:sources}
\begin{tabularx}{\textwidth}{lXl}
\toprule
Source & Content used & Layer of the network\\
\midrule
SET company pages & Top-10 shareholders at the latest book-closing date; sector; market value & Owners of listed firms\\
Form 56-1 One Report & Shareholding-structure section; notes on investments in subsidiaries, associates and joint ventures & Subsidiaries of listed firms\\
SEC Form 246-2 & Acquisitions and disposals crossing each 5\% threshold & Changes in large stakes\\
DBD (BOJ.5 and statements) & Shareholder lists and investment notes of private companies & Private holding companies\\
Bank Pillar 3 reports & Composition of the financial business group & Bank groups\\
GLEIF Level 2 & Direct and ultimate accounting parents of LEI holders & Cross-check\\
Licence registers & Banks, securities, asset-management and insurance licences & Financial-institution flags\\
\bottomrule
\end{tabularx}
\end{table}

\begin{table}[t]
\centering\small
\caption{Coverage of the network}\label{tab:coverage}
\input{generated/tab_coverage.tex}
\end{table}

\section{Measurement}\label{sec:method}
\paragraph{Network.} Let $w_{ij}\in(0,1]$ be the fraction of entity $j$'s shares held directly by entity $i$, and $M$
the matrix of these weights. Persons sharing a surname are aggregated into a family node, a reasonable default in
Thailand where surnames are registered and rarely shared by unrelated families; a manual mapping overrides it.

\paragraph{Cash-flow rights.} The share of $c$'s profits accruing to $a$ through all chains is the $(a,c)$ element of
\begin{equation}
\mathrm{CF} \;=\; \sum_{k\ge 1} M^{k} \;=\; M\,(I-M)^{-1},
\end{equation}
which remains well defined with cross-holdings \citep{brioschi1989}.

\paragraph{Control rights.} Following \citet{laporta1999} and \citet{claessens2000}, control along a chain is set by its
weakest link. We apply the rule block by block, so that one block of votes is never counted twice:
\begin{equation}
\mathrm{CR}(a,c)=\min\Big\{1,\;\sum_{s\,\in\,\mathcal{P}(c)} \min\big\{w_{sc},\,\mathrm{CR}(a,s)\big\}\Big\},\qquad \mathrm{CR}(a,a)=1,
\end{equation}
where $\mathcal{P}(c)$ is the set of direct shareholders of $c$. Nodes on ownership cycles (\NCycles{} cycles in the
data) are solved by fixed-point iteration.

\paragraph{Ultimate owners and groups.} In a first stage, the candidate owners are terminal nodes: families, persons,
state bodies, the Crown Property Bureau, foreign parents and private companies whose own shareholders are unknown.
Thai NVDR, custodians, funds and state investment funds (such as the Social Security Office) are excluded. A firm's
ultimate owner is the candidate with the largest control right at or above the threshold $t$ (\ThresholdPct{} in the
main specification; 10\% and 50\% as robustness). Firms without one are widely held; in a second stage, firms they
control are assigned to them, which separates ``widely held corporation'' and ``widely held financial'' apexes as in
\citet{laporta1999}. A group is the set of firms sharing an ultimate owner. For firms controlled by the state or the
Crown Property Bureau, groups are defined one level down, at the top firm (for example an energy or airport holding),
because merging all state enterprises into one group would hide their separate structures; ultimate-owner shares are
unaffected by this choice.

\paragraph{Pyramids and the wedge.} The depth of a firm is the length of the shortest control chain from its ultimate
owner along links controlled at $t$ or more. A firm sits in a pyramid when its depth is two or more. The wedge is
$\mathrm{CR}-\mathrm{CF}$ \citep{claessens2000,almeida2006}.

\paragraph{Finance--business nexus.} A bank group is linked to a business group through (i) ownership of at least 1\%
in either direction, (ii) a person on the board of a financial institution in the bank group and on a board in the
business group \citep{mizruchi1996}, or (iii) funds managed by the bank group's asset-management company holding
shares in the business group's listed firms.

\section{Results}\label{sec:results}
\subsection{The map}
Figure~\ref{fig:galaxy} draws every group with at least three firms. The largest, \TopGroupLabel, controls
\TopGroupFirms{} firms in \TopGroupIndustries{} SET industries, and the most diversified group spans
\MaxIndustries{} industries. \NGroupsMultiListed{} groups control two or more listed firms; together they account for
\ShareAffiliated{} of listed firms. The network is highly connected: the largest weakly connected component contains
\LargestWCCShare{} of all entities. Figure~\ref{fig:tree} shows the full control tree of the largest group.

\begin{figure}[p]
\centering
\includegraphics[width=\textwidth]{figures/galaxy_paper.png}
\caption{Thai business groups. Each circle is a group with at least three firms (area proportional to the number
of firms); the ultimate owner sits at the centre and controlled firms lie on rings by ownership distance. Ring colour
gives the type of ultimate owner, node colour the industry, and arcs the strongest links between groups (blue: shared
directors; orange: cross-ownership). Numbers match the ranks in Table~\ref{tab:groups}.}\label{fig:galaxy}
\end{figure}

\subsection{Who controls listed firms}
Table~\ref{tab:uo} and Figure~\ref{fig:uo} report ultimate owners. At the \ThresholdPct{} threshold, the most common
ultimate owner is the \DominantUOType{} category; families control \ShareFamilyTwenty, the state \ShareStateTwenty,
foreign parents \ShareForeignTwenty{} and the Crown Property Bureau \ShareCrownTwenty{} of listed firms, while
\ShareWidelyTwenty{} are widely held. A residual \ShareUnresolvedTwenty{} ends at private holding companies whose
shareholders are not yet in the data; they are the priority for further DBD requests. At the 10\% threshold, the
family share is \ShareFamilyTen.

\begin{table}[t]
\centering\small
\caption{Ultimate owners of listed firms (\% of firms; last column weighted by market value)}\label{tab:uo}
\input{generated/tab_uo.tex}
\end{table}

\begin{figure}[t]
\centering
\includegraphics[width=\textwidth]{figures/fig_uo.pdf}
\caption{Ultimate owners of listed firms by control threshold.}\label{fig:uo}
\end{figure}

\subsection{Size and concentration of groups}
Table~\ref{tab:groups} lists the largest groups. Group size is highly skewed (Figure~\ref{fig:groups}): the ten
largest groups by market value hold \TopTenMcapShare{} of the market value of listed firms, and the
Herfindahl--Hirschman index across groups is \HHIGroups.

\begin{table}[t]
\centering\footnotesize
\caption{The fifteen largest groups by number of firms}\label{tab:groups}
\input{generated/tab_groups.tex}
\end{table}

\begin{figure}[t]
\centering
\includegraphics[width=\textwidth]{figures/fig_groups.pdf}
\caption{Distribution of group size and depth.}\label{fig:groups}
\end{figure}

\subsection{Pyramids and the separation of ownership and control}
Among listed firms with an ultimate owner, \SharePyramid{} are held through pyramids. The mean wedge between control
and cash-flow rights is \MeanWedge{} percentage points (median \MedianWedge), and the mean ratio of cash-flow to control
rights is \MeanCFCRRatio. Control chains reach \MaxDepth{} links; the average firm in the network sits \MeanDepth{}
links below its ultimate owner (Figure~\ref{fig:wedge}). \NFirmsInCycles{} firms belong to \NCycles{} ownership cycles.

\begin{figure}[t]
\centering
\includegraphics[width=\textwidth]{figures/fig_wedge.pdf}
\caption{Control versus cash-flow rights of listed firms, and the length of control chains.}\label{fig:wedge}
\end{figure}

\subsection{The finance--business nexus}
The data contain \NFIEntities{} financial institutions, including \NBanks{} banks, organised in \NBankGroups{} bank
groups. We find \NNexusPairs{} linked pairs of bank groups and business groups: \NDirectorLinks{} director links,
\NOwnershipLinks{} ownership links and \NFundLinks{} fund links. Bank groups share at least one director with
\ShareBigGroupsBankInterlock{} of business groups that have ten or more firms (Figure~\ref{fig:nexus},
Table~\ref{tab:nexus}). Fund links reflect portfolio investment by asset-management arms on behalf of unit holders,
not control by the bank.

\begin{figure}[t]
\centering
\includegraphics[width=\textwidth]{figures/fig_nexus.pdf}
\caption{Links between bank groups and the most connected business groups.}\label{fig:nexus}
\end{figure}

\begin{table}[t]
\centering\footnotesize
\caption{Strongest bank--business group links}\label{tab:nexus}
\input{generated/tab_nexus.tex}
\end{table}

\subsection{Robustness and validation}
Table~\ref{tab:robust} recomputes ultimate owners at alternative thresholds, without family aggregation, and treating
state investment funds as owners. GLEIF reports \GLEIFPairs{} parent relationships for Thai entities in our data; the
network agrees with \GLEIFAgreement{} of them (same group or an ownership path from parent to child).
\ValidationStatement

\begin{table}[t]
\centering\small
\caption{Robustness: ultimate owners of listed firms under alternative specifications (\% of firms)}\label{tab:robust}
\input{generated/tab_robust.tex}
\end{table}

\begin{figure}[p]
\centering
\includegraphics[width=\textwidth]{figures/tree_top.pdf}
\caption{Control tree of the largest group. The ultimate owner is at the centre; each ring is one more ownership link.
Listed companies are labelled with ticker and name; squares mark holding
companies.}\label{fig:tree}
\end{figure}

\section{Discussion and limitations}\label{sec:discussion}
Four limitations bound the interpretation. First, chains that leave Thailand through holding companies registered
abroad stop at the border unless the foreign parent is itself disclosed, so foreign ultimate owners may be
intermediate vehicles. Second, nominee shareholders and voting agreements are invisible in share registers; control
can be stronger, or held by different people, than the percentages suggest. Third, surname-based family aggregation
can merge unrelated people who share a surname or split families whose members use different surnames; the manual
mapping and the robustness table address this only in part. Fourth, sources differ in their reference dates (book
closing, fiscal year end, annual general meeting), so the network is a snapshot assembled from several dates within
one year. Accounting control under TFRS~10 can also differ from the share-based definition used here.

Within these limits, the map supports descriptive conclusions about the reach of groups, the use of pyramids and the
ties between banks and business groups. Causal questions, such as whether pyramids are associated with lower minority
shareholder value \citep{bertrand2002,masulis2011}, require panel data across years and are left for future work.

\section{Conclusion}
Built entirely from public disclosures, this map shows how Thai business groups extend from a few ultimate owners
through listed firms into thousands of unlisted subsidiaries, and how banks connect to them. Because every step of the
pipeline is scripted, the map can be rebuilt each year and extended, for example to study changes in group structure,
related-party transactions or market concentration.

\bibliographystyle{plainnat}
\bibliography{references}

\appendix
\section{Reproducing the results}
\begin{small}
\begin{verbatim}
pip install -r requirements.txt
python -m tbg init                      # create folders and blank input files
python -m tbg collect set-universe      # or place the SET list in data/raw/universe/
python -m tbg extract pages --kind subsidiaries
python -m tbg extract run   --kind subsidiaries --batch
python -m tbg extract collect-batch --kind subsidiaries
python -m tbg collect gleif
python -m tbg all                       # build, analyze, figures, explorer, paper
\end{verbatim}
\end{small}
\end{document}
